% ECONOMICS_PROBLEM: {"id": "D83-521", "title": "Optimal information disclosure", "jel_code": "D83", "source_id": "OP-0402"}
% !TeX program = xelatex
\documentclass[11pt,a4paper]{article}
\usepackage[margin=23mm]{geometry}
\usepackage{fontspec}
\setmainfont{Latin Modern Roman}
\usepackage{amsmath,amssymb,mathtools}
\usepackage{xcolor,enumitem,booktabs,longtable,array,xurl,hyperref}
\definecolor{muted}{HTML}{53636C}
\hypersetup{colorlinks=true,urlcolor=blue,linkcolor=blue}
\setlength{\parindent}{0pt}
\setlength{\parskip}{5pt}
\newcommand{\R}{\mathbb R}
\newcommand{\E}{\mathbb E}
\newcommand{\N}{\mathbb N}
\newcommand{\ind}{\mathbf 1}
\newcommand{\Var}{\operatorname{Var}}
\newcommand{\Cov}{\operatorname{Cov}}
\newcommand{\supp}{\operatorname{supp}}
\DeclareMathOperator*{\argmin}{arg\,min}
\DeclareMathOperator*{\argmax}{arg\,max}
\newcommand{\entrymeta}[3]{{\sffamily\small\color{muted}\textbf{\texttt{#1}}\quad|\quad #2\par #3\par}}
\newcommand{\Problemref}[1]{\texttt{#1}}
\newcommand{\JELref}[2]{\texttt{#2} (JEL \texttt{#1})}
\begin{document}
\section*{Optimal information disclosure}
\textbf{Problem ID:} \texttt{D83-521}\quad \textbf{JEL:} \texttt{D83}\par
% BEGIN ECONOMICS STATEMENT
\label{problem:OP-0402}
{\sffamily\small\color{muted}Original subject group: Digital, platform and data economics\par}
\label{definitions:problem:30007684}
\textbf{Audit classification:} Published research agenda. The attention/disclosure question is verified. The original unrestricted argmax need not exist and finite known welfare ranking is not inherently open.
\subsection*{Definitions and precise research target}
Consider a service marketplace whose buyers must compare sellers under a fixed attention budget. Each seller has a characteristic vector \(x\in\mathbb R^K\), where the positive integer \(K\) is the number of possible disclosed attributes; these include quality, price and return rates. A disclosure rule \(d\) maps verified characteristics and seller reports to a displayed message. A feasible rule discloses at most \(B\) attributes, where \(1\leq B<K\), and specifies which attributes must use standardized units. Buyers choose whether to search further and which seller to hire, while sellers may change quality and voluntarily disclose allowed information. Let \(\mathcal D_B\) be a specified finite menu of implementable rules satisfying the attention limit, truthful verification constraints and prespecified minimum safety disclosures. Assume these requirements are jointly feasible, so \(\mathcal D_B\neq\varnothing\). Define \(W(d)\) as expected consumer surplus plus supplier and platform profits after these responses, with profits net of production and verification costs and surplus net of attention costs, for the platform's registered-buyer population during one year. The design target is \(d^*\in\arg\max_{d\in\mathcal D_B}W(d)\). Establish sufficient observable restrictions or experimental comparisons to rank candidate rules, including standardized disclosure versus seller discretion. If buyer attention or seller responses are only partially identified, report a welfare range for each rule. This provides a finite-budget design agenda rather than assuming that releasing more information always improves choice.

\textbf{Definitions, assumptions and mathematical qualifications.} For an existence-safe formulation take a finite nonempty menu \(\mathcal D_B\) of implementable disclosure rules. Each rule specifies which messages are shown after verified measurements and reports, their standardized units, allowed seller withholding, buyer inspection costs and mandatory safety content. Define feasible consumer search and seller-response equilibria with a selection kernel. Welfare is measured under one population and common monetary units, including verification and attention resource costs. Then \(\arg\max_{d\in\mathcal D_B}W(d)\) is nonempty when welfare is finite. An unrestricted measurable rule class instead needs compactness and upper semicontinuity or must use \(\sup_dW(d)\) and \(\varepsilon\)-optimal rules.
\subsection*{Variants and assumption changes}
\begin{enumerate}
\item \textbf{Editorial, finite known menu.} If each \(W(d)\) is known, select its maximizer by finite comparison; this is a solved optimization benchmark.
\item \textbf{Editorial, unidentified responses.} Observe only a subset of disclosure experiments; maximize worst-case welfare over compatible buyer attention and seller-response models, and choose new experiments to reduce maximum policy regret.
\end{enumerate}
\subsection*{References and source audit}
\textbf{Checked source.} 2025 conference draft, Section 1.1, second research direction, limited-attention paragraph, PDF p. 5. Full primary text retrieved. \url{https://conference.nber.org/confer/2025/DTs25/farronato.pdf}.
\begin{enumerate}\raggedright
\item\hypertarget{definitions:source:30007684:1}{} Chiara Farronato. 2025. \emph{Digital Platforms as Information Aggregators: Implications for their Design and Regulation}. 2025 conference draft, Section 1.1, second research direction, limited-attention paragraph, PDF p. 5.
\url{https://conference.nber.org/confer/2025/DTs25/farronato.pdf}.
\end{enumerate}

\subsection*{Common conventions: Source mathematical conventions}

These conventions apply to each entry unless explicitly replaced.

\textbf{Models and identification.} A primitive model \(m\) specifies all probability laws, feasible choices, information, equilibrium and measurement rules used by its target. The admissible class \(\mathcal M\) is fixed independently of outcomes. If \(P_O(m)\) is its observable probability law and \(\tau(m)\) the target, its identified set at observable law \(P\) is
\[
 \mathcal I_\tau(P)=\{\tau(m):m\in\mathcal M,\ P_O(m)=P\}.
\]
Point identification means that this set is a singleton. An empty set signals incompatibility of data and assumptions. Coordinate bounds need not describe the joint set. Bounds are sharp only if every retained target is attainable or arbitrarily approachable by compatible models. Finite bounds require explicit restrictions on unobserved quantities; unrestricted confounding, hidden wealth, extrapolation or arbitrary equilibrium selection can yield uninformative sets.

\textbf{Causal interventions.} A potential outcome \(Y_i(p)\) is the outcome under a complete policy regime \(p\), including timing, financing, information and market spillovers. The target population and units are fixed before treatment. With interference, use the complete assignment vector or a justified exposure mapping; individual no-interference notation alone cannot identify an economy-wide intervention. A difference of mean potential outcomes does not require the cross-world joint distribution. Conditioning on post-treatment migration, survival, enrollment or employment changes the estimand and needs separate justification.

\textbf{Statistical statements.} An estimator, confidence set, forecast or test specifies a sampling law, model class and loss/coverage criterion. Uniform \((1-\alpha)\)-coverage means \(\inf_{m\in\mathcal M}\Pr_m\{\tau(m)\in C_n\}\geq1-\alpha\), or an explicitly asymptotic version. The whole parameter space trivially has coverage, so informativeness requires a stated width, risk or power objective. A forecasting improvement is relative to a named benchmark, information set and evaluation population.

\textbf{Equilibrium and optimization.} An equilibrium comprises feasible plans, optimality, beliefs where needed, budget constraints and all market/resource clearing. The primitives or a stated benchmark define each correspondence \(\mathcal E(p;m)\). Nonemptiness is assumed or proved, never inferred just from compact policy sets. A selected-equilibrium target uses a declared selection rule; a robust target quantifies over all permitted equilibria. Use suprema or infima unless attainment is established. An \(\varepsilon\)-optimal policy is feasible and within \(\varepsilon\) of the finite optimum value. Report an empty feasible set separately.

\textbf{Welfare and accounting.} Normative utility, interpersonal normalization, social weights, numeraire, discounting and the use of public receipts are primitives. Behavioral utility need not coincide with the welfare criterion. Household welfare includes ownership income and rents once; adding profits again to compensating variation generally double-counts them. A monetary equivalent is defined by an explicit transfer and price convention, with monotonicity and range conditions. Average effects, mechanism contributions, transfers and welfare losses are different targets. Infinite sums require convergence and dynamic feasible plans require terminal or no-Ponzi conditions.

\textbf{Source versions.} A working-paper date, revision date and journal publication year are distinguished. A DOI for a journal version is not automatically the DOI of an earlier manuscript. Locators refer to the stated version and printed pages where specified. A metadata-only check does not verify a claim about a full-text conclusion. Every entry's source subsection narrows the provenance claim accordingly.
% END ECONOMICS STATEMENT
\end{document}
